\subsection{Artin 阴影的常数层压缩与同余偏置的 finite-part 定域化}

正文小节 \ref{subsec:conclusion-realinput40-primitive-surgery-residue-escort-geometry}
以及附录真实输入 \(40\) 状态核的 primitive/arity 分析，已经分别给出
\(\zeta_M\)
的三因子分解、主极点 \(z_\star=\varphi^{-2}\) 处的有限部加法律，以及 arity 模类的 Abel--Mertens 常数。
把这两条链拼接后，可把“Artin 阴影”与“同余偏置”同时压缩为常数层现象。

\begin{theorem}[Artin 阴影的常数层压缩与极点唯一承载]\label{thm:derived-artin-shadow-constant-layer-compression}
\leanverified{paper\_derived\_artin\_shadow\_constant\_layer\_compression}
设
\[
\zeta_M(z)=\zeta_{\mathrm{triv}}(z)\,\zeta_{\mathrm{even}}(z)\,\zeta_{\chi}(z),
\]
其中
\[
\zeta_{\mathrm{triv}}(z)=\frac{1}{(1-z)^2(1+z)^3},
\qquad
\zeta_{\mathrm{even}}(z)=\frac{1}{1-3z+z^2},
\qquad
\zeta_{\chi}(z)=\frac{1}{1+z-z^2}.
\]
记主极点
\[
z_\star=\varphi^{-2}.
\]
则在 \(z=z_\star\) 处，仅 \(\zeta_{\mathrm{even}}\) 具有极点，而
\(
\zeta_{\mathrm{triv}}
\)
与
\(
\zeta_\chi
\)
都解析且不为零。
因此主奇性完全由 even 通道承担，而 trivial/sign 两个 Artin 阴影只在常数层留下有限偏移：
\[
\log\mathfrak{M}
=
\log\mathfrak{M}_{\mathrm{even}}
+K_{-F}
+P_{\mathrm{triv}}(z_\star),
\qquad
P_{\mathrm{triv}}(z_\star)=9-4\sqrt5.
\]
换言之，真实输入 \(40\) 状态核的主极点几何是纯 even 的；所有非平凡 Artin 贡献都被压缩为 finite-part 常数修正。
\end{theorem}

\begin{proof}
由注 \ref{rem:real-input-40-logM-artin-factorization}，
\(
\zeta_M
\)
确有上述三因子分解。
又
\[
1-3z_\star+z_\star^2=0,
\]
故 \(\zeta_{\mathrm{even}}\) 在 \(z_\star\) 处有极点。
另一方面，
\[
(1-z_\star)^2(1+z_\star)^3\neq 0,
\qquad
1+z_\star-z_\star^2\neq 0,
\]
故 \(\zeta_{\mathrm{triv}}\) 与 \(\zeta_\chi\) 在 \(z_\star\) 处均解析且非零。
于是主极点的留数与主增长项都只能由 even 通道承载。

同一注已经给出常数层分解
\[
\log\mathfrak{M}
=
\log\mathfrak{M}_{\mathrm{even}}+K_{-F}+P_{\mathrm{triv}}(z_\star).
\]
再由命题 \ref{prop:triv-factor-primitive-polynomial}，
\[
P_{\mathrm{triv}}(z)=-z+3z^2.
\]
代入 \(z=z_\star=\varphi^{-2}\) 即得
\[
P_{\mathrm{triv}}(z_\star)
=
-\varphi^{-2}+3\varphi^{-4}
=
9-4\sqrt5.
\]
证毕。
\end{proof}

\begin{theorem}[同余偏置的常数层定域化]\label{thm:derived-congruence-bias-finitepart-localization}
\leanverified{paper\_derived\_congruence\_bias\_finitepart\_localization}
对真实输入 \(40\) 状态核的 arity 模 \(m\) 分类，记
\[
N_{n,r}^{(m)}
\]
为 residue 类 \(r\in \ZZ/m\ZZ\) 的 primitive 计数，\(p_n\) 为总 primitive 计数，\(\lambda=\varphi^2\) 为 Perron 根。
则对每个 \(r\) 都有统一渐近
\[
\sum_{n\le N}\frac{N_{n,r}^{(m)}}{\lambda^n}
=
\frac{1}{m}\log N + C_{m,r}+o(1),
\qquad
C_{m,r}=\mathsf{M}_r^{(m)}.
\]
因而任意两类 \(r,s\in\ZZ/m\ZZ\) 之间的差异完全下沉为 finite-part 常数：
\[
\sum_{n\le N}\frac{N_{n,r}^{(m)}-N_{n,s}^{(m)}}{\lambda^n}
=
\mathsf{M}_r^{(m)}-\mathsf{M}_s^{(m)}+o(1).
\]
特别地，在 \(m=2\) 时，Chebyshev 型偏置
\[
B_n:=N_{n,0}^{(2)}-N_{n,1}^{(2)}
\]
满足
\[
\frac{B_n}{p_n}
=
O\!\left(\left(\frac{\rho(M(-1))}{\lambda}\right)^n\right),
\qquad
\frac{\rho(M(-1))}{\lambda}\approx 0.708038612.
\]
因此，同余偏置既不会改变主极点位置，也不会扰动 Perron 标度；它完全是 Abel/Mertens 有限部层上的定域现象。
\end{theorem}

\begin{proof}
附录 arity 分析已经证明：对每个 residue 类 \(r\)，都有
\[
\sum_{n\le N}\frac{N_{n,r}^{(m)}}{\lambda^n}
=
\frac{1}{m}\log N + C_{m,r}+o(1),
\qquad
C_{m,r}=\mathsf{M}_r^{(m)}.
\]
于是主项 \(\frac{1}{m}\log N\) 与 \(r\) 无关。
对两类 \(r,s\) 相减，主项完全抵消，得到
\[
\sum_{n\le N}\frac{N_{n,r}^{(m)}-N_{n,s}^{(m)}}{\lambda^n}
=
\mathsf{M}_r^{(m)}-\mathsf{M}_s^{(m)}+o(1),
\]
这正表明类间偏置全部定域于 finite-part 常数。

当 \(m=2\) 时，同一附录进一步给出
\[
B_n=O\!\left(\frac{\rho(M(-1))^n}{n}\right).
\]
再结合真实输入 \(40\) 状态核的 primitive 主项
\(
p_n\asymp \lambda^n/n
\)，便得到
\[
\frac{B_n}{p_n}
=
O\!\left(\left(\frac{\rho(M(-1))}{\lambda}\right)^n\right),
\]
其中数值比值约为 \(0.708038612<1\)。
故二元 parity 偏置相对于总 primitive 规模呈指数小，进一步说明同余信息不会进入主奇性层。证毕。
\end{proof}

\begin{remark}[主奇性与类偏置的严格分层]\label{rem:derived-artin-even-finitepart-congruence-splitting}
定理 \ref{thm:derived-artin-shadow-constant-layer-compression}
与定理 \ref{thm:derived-congruence-bias-finitepart-localization}
共同说明：在真实输入 \(40\) 状态核上，主极点几何、Artin 阴影与同余类偏置并不处于同一层。
主奇性由 even 因子独占；sign/trivial 通道只在常数层留下修正；而模 \(m\) 的 residue 偏置则更进一步，仅表现为逐类 Abel--Mertens 常数之差。
\end{remark}

\endinput
